\section{Discussion and Open Problems}
\label{sec:discussion}

The central contribution of this paper is to expose a semantic mismatch
between cross-validation diagnostics and adversarial replacement robustness
for deterministic local learning rules. Leave-one-out stability is a statement
about the effect of deleting a point and evaluating at that same point;
replace-one stability concerns the effect of replacing a point by an
adversarially chosen labeled example and evaluating at any query. The paper
shows that these two notions are structurally distinct, through a systematic
perturbation calculus, an exact margin-crossing criterion, and a general
separation theorem for $k$-NN that holds for every $k \ge 1$.

\subsection{Relationship to Other Notions of Stability}

It is useful to situate the present work relative to three established lines
of inquiry that also use the term ``stability.''

\paragraph{Generalization stability (Bousquet--Elisseeff).}
Classical algorithmic stability \cite{bousquet2002stability} studies how
replacing one training example affects the expected or fresh-test loss, with
the goal of bounding generalization error. That agenda is distributional and
asymptotic. Our work is combinatorial and worst-case: we fix a finite sample
and ask whether a single perturbation changes the prediction at a specific
query. There is no contradiction between a rule being stable in the
generalization sense and having the pointwise LOO/replace-one gap we
demonstrate.

\paragraph{Consistency (Cover--Hart, Stone).}
Classical $k$-NN theory \cite{cover1967nearest,stone1977consistent} studies
whether the classifier converges to the Bayes risk as $n \to \infty$ under
distributional assumptions. Those are limit theorems about expected
performance. Our witness constructions are finite and adversarial. A rule can
be consistent and simultaneously admit finite configurations where LOO and
replace-one diverge.

\paragraph{Conformal prediction stability.}
Recent work \cite{liang2023algorithmic,pournaderi2024training} studies how
deleting or exchanging points affects prediction sets and coverage
certificates. Those analyses operate under exchangeability or
training-conditional assumptions and typically consider real-valued conformity
scores rather than 0--1 loss. Our results are complementary: they show that
even in the simpler 0--1 loss setting, the choice of perturbation operation
matters before any distributional or coverage considerations enter.

\subsection{Summary of Established Results}

\begin{enumerate}[leftmargin=*,label=(\arabic*)]
    \item \textbf{Uniform decomposition:} replace-one $\le$ delete-one +
        insert-one (\Cref{prop:4.1}); the decomposition is algebraically sharp.
    \item \textbf{Margin criterion:} deterministic $k$-NN prediction changes
        iff the signed vote margin crosses zero (\Cref{prop:5.1} and
        \Cref{prop:generalized-margin} for DLVRs).
    \item \textbf{General separation:} for every odd $k \ge 3$, there exists a
        finite metric sample with $\Delta_{\mathrm{loo}}(S,i)=0$ for all $i$
        and $\Delta_{\mathrm{rep}}(S,j,z,x,y)=1$
        (\Cref{thm:odd-k-separation}).
    \item \textbf{Even-$k$ separation:} the analogous result holds for even
        $k \ge 2$ under deterministic tie-breaking
        (\Cref{thm:even-k-separation}), with a tighter dependence on the
        tie-breaking convention.
    \item \textbf{Stability hierarchy:} the relationships among stability
        notion pairs are mapped with theorem, counterexample, or open-question
        status (\Cref{sec:stability-hierarchy}).
    \item \textbf{Abstract framework:} the uniform calculus and margin criterion
        extend to deterministic local vote rules
        (\Cref{sec:local-vote-framework}).
\end{enumerate}

\subsection{Practical Implications}

The LOO/replace-one gap has concrete consequences for practitioners using
$k$-NN or similar local rules. Low LOO error is often interpreted as evidence
of model stability, but a model can be perfectly LOO-stable while being highly
vulnerable to a single adversarially chosen replacement. The experiments
(\Cref{sec:tabular-benchmark}) confirm that the gap appears on standard UCI
datasets at measurable frequency. The margin-based diagnostic (line~7 of
\Cref{alg:diagnostic}) provides a computationally efficient way to detect
vulnerable queries: queries with $|M_k(S,x)| \le 2$ are potentially vulnerable
to a single label-flip replacement. Practitioners should therefore not equate
LOO stability with replacement robustness, and should monitor the margin
distribution rather than relying solely on aggregate LOO error.

\subsection{Limitations}

\begin{enumerate}[leftmargin=*,label=(\arabic*)]
    \item \textbf{Beyond binary 0--1 loss.} The calculus is developed for
        binary classification with 0--1 loss. Extending it to real-valued
        scores, multiclass labels, or regression loss functions would require
        a more general notion of margin crossing.
    \item \textbf{Beyond finite metric spaces.} The constructions use finite
        metric spaces (often graph metrics). Generalizing to infinite spaces,
        continuous densities, or non-metric similarity measures is open.
    \item \textbf{Beyond $k$-NN.} The abstract local vote rule framework
        (\Cref{sec:local-vote-framework}) provides the vocabulary, but the
        explicit separation constructions are specific to $k$-NN. Finding
        analogous constructions for kernel rules, radius-neighbor classifiers,
        or decision trees would strengthen the claim of generality.
    \item \textbf{Randomized tie-breaking.} Our analysis is for deterministic
        tie-breaking. The even-$k$ results show that randomization can close
        some gaps. Characterizing the stability hierarchy under randomized
        tie-breaking is an open problem.
    \item \textbf{Distributional connections.} The paper studies worst-case
        fixed-sample stability. Connecting the finite combinatorial guarantees
        to distributional stability notions (expected replace-one,
        high-probability bounds) without collapsing the distinctions we
        isolate is an important next step.
\end{enumerate}

\subsection{Open Problems}

\begin{enumerate}[leftmargin=*,label=(\arabic*)]

    \item \textbf{Tie-free even-$k$ separation.}
    Does there exist, for any even $k \ge 2$, a finite metric sample for
    deterministic $k$-NN in which the original prediction at the designated
    query is determined by a strict majority (margin $\neq 0$) and the
    LOO/replace-one separation holds? If not, the even-$k$ separation is
    inherently tied to tie-breaking.

    \item \textbf{Characterization of DLVRs admitting separation.}
    Characterize the class of deterministic local vote rules for which the
    LOO/replace-one separation holds. Is it exactly the class of rules whose
    margin function can change by at least $2 \cdot \min_i w(x,x_i)$ under a
    single replacement while remaining invariant under any single deletion at
    the evaluation point?

    \item \textbf{Quantitative bounds.}
    For a given data distribution, what is the expected LOO/replace-one gap?
    Can the worst-case gap be bounded by a function of the data geometry (e.g.,
    the doubling dimension or the number of duplicate-prone points)?

    \item \textbf{Algorithmic mitigation.}
    Can the StabilityGapDiagnostic be extended to an algorithm that
    \emph{repairs} vulnerable queries, e.g., by recommending a different $k$,
    a different tie-breaking rule, or a specific subset of points to re-label?

    \item \textbf{Connection to adversarial robustness.}
    The replace-one perturbation is a form of training-set adversarial example.
    How does the LOO/replace-one gap relate to standard test-time adversarial
    examples? Is there a transfer between training-set and test-set
    vulnerabilities?

\end{enumerate}

\subsection{Closing Reflection}

If there is a broader moral, it is this: for local learning rules, the geometry
of the neighborhood and the semantics of the perturbation must be specified
together. Once that is done, some relationships become cleaner than expected
(the uniform calculus), and some purportedly similar notions separate much
sooner than intuition suggests (the LOO/replace-one gap). The scale at which
these separations appear -- already at two metric points and $k=1$ -- suggests
that the phenomenon is not a high-complexity artifact but a basic feature of
local decision rules.
